%% ---------------------------------------------------------------------
%%  FIGURA - Pipeline SDD com portoes de verificacao (DG-03)
%% ---------------------------------------------------------------------
\begin{figure}[htb]
	\caption{Pipeline SDD com os seis portões de verificação}
	\label{fig:pipeline-gates}
	\centering
	\begin{tikzpicture}[
			font=\footnotesize,
			etapa/.style={draw, rounded corners=2pt, minimum width=1.95cm, minimum height=0.9cm, align=center, fill=black!4},
			portao/.style={draw, fill=black!12, rounded corners=6pt, minimum width=0.85cm, minimum height=0.5cm, font=\scriptsize, align=center},
			>=Stealth,
			fl/.style={->, thick}
		]

		\node[etapa] (e1) at (0,0) {Estado\\ do repositório};
		\node[portao, right=0.35cm of e1] (g0) {G0};
		\node[etapa, right=0.35cm of g0] (e2) {Intenção\\ e requisito};
		\node[portao, right=0.35cm of e2] (g1) {G1};
		\node[etapa, right=0.35cm of g1] (e3) {Projeto e\\ decisões};
		\node[portao, right=0.35cm of e3] (g2) {G2};
		\node[etapa, below=1.1cm of e3] (e4) {Implementação\\ mínima};
		\node[portao, left=0.35cm of e4] (g3) {G3};
		\node[etapa, left=0.35cm of g3] (e5) {Evidência:\\ build e testes};
		\node[portao, left=0.35cm of e5] (g4) {G4};
		\node[etapa, left=0.35cm of g4] (e6) {Encerramento\\ e estado};
		\node[portao, left=0.35cm of e6] (g5) {G5};

		\draw[fl] (e1) -- (g0);
		\draw[fl] (g0) -- (e2);
		\draw[fl] (e2) -- (g1);
		\draw[fl] (g1) -- (e3);
		\draw[fl] (e3) -- (g2);
		\draw[fl] (g2) -- (e4);
		\draw[fl] (e4) -- (g3);
		\draw[fl] (g3) -- (e5);
		\draw[fl] (e5) -- (g4);
		\draw[fl] (g4) -- (e6);
		\draw[fl] (e6) -- (g5);
	\end{tikzpicture}
	\legend{Fonte: elaborado pelo autor (2026) a partir de \texttt{docs/agentic/WORKFLOW.md}.
	O portão G2 é aplicado apenas quando a complexidade exige decisão arquitetural.}
\end{figure}
